\documentclass[11pt,a4paper]{article}

% ====== PAQUETES Y CONFIGURACIÓN ======
\usepackage[utf8]{inputenc}
\usepackage[spanish]{babel}
\usepackage{geometry}
\geometry{a4paper, margin=2.5cm}
\usepackage{amsmath,amsfonts,amssymb}
\usepackage{graphicx}
\usepackage{xcolor}
\usepackage{pagecolor}
\usepackage{afterpage}
\usepackage{tcolorbox}
\tcbuselibrary{listings}
\usepackage{listings}
\usepackage{fancyhdr}
\usepackage{titlesec}
\usepackage{hyperref}
\usepackage{fontawesome5}
\usepackage{enumitem}
\usepackage{multicol}
\usepackage{tikz}
\usepackage{pgfplots}
\pgfplotsset{compat=1.18}

% ====== DEFINICIÓN DE COLORES K8S OSCUROS ======
\definecolor{k8sblue}{RGB}{0,123,255}           % Azul Kubernetes oficial
\definecolor{k8sdarkblue}{RGB}{13,71,161}       % Azul oscuro principal
\definecolor{k8snavyblue}{RGB}{25,46,94}        % Azul marino muy oscuro
\definecolor{k8slightblue}{RGB}{63,149,244}     % Azul claro para acentos
\definecolor{k8scyan}{RGB}{0,188,212}           % Cyan para información
\definecolor{k8ssteelblue}{RGB}{69,90,120}      % Azul acero para secundarios
\definecolor{k8scodebg}{RGB}{22,29,47}          % Fondo código azul muy oscuro
\definecolor{k8swarning}{RGB}{255,152,0}        % Naranja para warnings
\definecolor{k8serror}{RGB}{244,67,54}          % Rojo para errores
\definecolor{k8ssuccess}{RGB}{76,175,80}        % Verde para success
\definecolor{k8slightgray}{RGB}{236,239,244}    % Gris claro
\definecolor{k8sdarkgray}{RGB}{96,108,118}      % Gris oscuro

% ====== CONFIGURACIÓN DE HYPERLINKS ======
\hypersetup{
    colorlinks=true,
    linkcolor=k8slightblue,
    urlcolor=k8scyan,
    citecolor=k8ssteelblue,
    pdfborderstyle={/S/U/W 1},
    pdftitle={Kubernetes CKA 2024 - Guía Completa},
    pdfauthor={Administrador Kubernetes},
    pdfsubject={Certificación CKA},
    pdfkeywords={Kubernetes, CKA, Administración, Containers}
}

% ====== CONFIGURACIÓN DE HEADERS CON LOGO K8S ======
\pagestyle{fancy}
\fancyhf{}
% Header izquierdo con logo K8s simulado y título
\fancyhead[L]{%
    \textcolor{k8sdarkblue}{\textbf{⎈}} \hspace{2pt}
    \textcolor{k8sdarkblue}{\textbf{Kubernetes CKA 2024}}
    \textcolor{k8ssteelblue}{\textbf{ - Guía Completa}}
}
\fancyhead[R]{\textcolor{k8sdarkgray}{\textbf{\thepage}}}
\renewcommand{\headrulewidth}{3pt}
\renewcommand{\headrule}{%
    \hbox to\headwidth{%
        \color{k8sdarkblue}\leaders\hrule height 2pt\hfill%
        \kern2pt%
        \color{k8slightblue}\leaders\hrule height 1pt\hfill%
    }%
}
\setlength{\headheight}{25pt}

% ====== ESTILOS DE TÍTULOS CON GRADIENTE AZUL ======
\titleformat{\section}
{\color{k8sdarkblue}\LARGE\bfseries}
{\textcolor{k8slightblue}{\thesection}}{1em}{}

\titleformat{\subsection}
{\color{k8snavyblue}\Large\bfseries}
{\textcolor{k8scyan}{\thesubsection}}{1em}{}

\titleformat{\subsubsection}
{\color{k8ssteelblue}\large\bfseries}
{\textcolor{k8slightblue}{\thesubsubsection}}{1em}{}

% ====== CAJAS PERSONALIZADAS CON TEMA AZUL OSCURO ======
\tcbset{
    commandstyle/.style={
        colback=k8scodebg,
        coltext=white,
        colframe=k8slightblue,
        boxrule=2.5pt,
        arc=8pt,
        left=12pt,
        right=12pt,
        top=10pt,
        bottom=10pt,
        shadow={2pt}{2pt}{0pt}{k8snavyblue!50},
        borderline west={4pt}{0pt}{k8scyan}
    }
}

\newtcolorbox{commandbox}{
    commandstyle,
    fontupper=\ttfamily\small,
    before upper={\parindent0pt}
}

\newtcolorbox{warningbox}{
    colback=k8swarning!5,
    coltext=black,
    colframe=k8swarning,
    boxrule=2pt,
    arc=6pt,
    left=12pt,
    right=12pt,
    top=10pt,
    bottom=10pt,
    title={\textcolor{k8swarning}{\faExclamationTriangle}\ \textcolor{k8sdarkblue}{\textbf{Importante}}},
    fonttitle=\bfseries,
    toptitle=8pt,
    bottomtitle=8pt,
    shadow={1pt}{1pt}{0pt}{k8swarning!30}
}

\newtcolorbox{infobox}{
    colback=k8scyan!8,
    coltext=black,
    colframe=k8scyan,
    boxrule=2pt,
    arc=6pt,
    left=12pt,
    right=12pt,
    top=10pt,
    bottom=10pt,
    title={\textcolor{k8scyan}{\faInfoCircle}\ \textcolor{k8sdarkblue}{\textbf{Información}}},
    fonttitle=\bfseries,
    toptitle=8pt,
    bottomtitle=8pt,
    shadow={1pt}{1pt}{0pt}{k8scyan!30}
}

\newtcolorbox{tipbox}{
    colback=k8ssuccess!8,
    coltext=black,
    colframe=k8ssuccess,
    boxrule=2pt,
    arc=6pt,
    left=12pt,
    right=12pt,
    top=10pt,
    bottom=10pt,
    title={\textcolor{k8ssuccess}{\faLightbulb}\ \textcolor{k8sdarkblue}{\textbf{Tip Profesional}}},
    fonttitle=\bfseries,
    toptitle=8pt,
    bottomtitle=8pt,
    shadow={1pt}{1pt}{0pt}{k8ssuccess!30}
}

\newtcolorbox{errorbox}{
    colback=k8serror!8,
    coltext=black,
    colframe=k8serror,
    boxrule=2pt,
    arc=6pt,
    left=12pt,
    right=12pt,
    top=10pt,
    bottom=10pt,
    title={\textcolor{k8serror}{\faTimes}\ \textcolor{k8sdarkblue}{\textbf{Error Común}}},
    fonttitle=\bfseries,
    toptitle=8pt,
    bottomtitle=8pt,
    shadow={1pt}{1pt}{0pt}{k8serror!30}
}

% ====== CONFIGURACIÓN DE LISTINGS CON TEMA OSCURO ======
\lstset{
    backgroundcolor=\color{k8scodebg},
    basicstyle=\color{white}\ttfamily\small,
    keywordstyle=\color{k8slightblue}\bfseries,
    stringstyle=\color{k8scyan},
    commentstyle=\color{k8sdarkgray}\itshape,
    numbers=left,
    numberstyle=\tiny\color{k8ssteelblue},
    stepnumber=1,
    numbersep=8pt,
    frame=leftline,
    rulecolor=\color{k8slightblue},
    framerule=3pt,
    xleftmargin=20pt,
    xrightmargin=10pt,
    breaklines=true,
    breakatwhitespace=false,
    showstringspaces=false,
    tabsize=4,
    captionpos=t,
    escapeinside={(*@}{@*)}
}

% ====== DOCUMENTO ======
\begin{document}

% ====== PORTADA PROFESIONAL CON DISEÑO K8S ======
\begin{titlepage}
    \pagecolor{k8snavyblue}
    \centering
    \vspace*{1.5cm}
    
    % Logo personalizado de Kubernetes con TikZ
    \begin{tikzpicture}[scale=1.2]
        % Círculo exterior azul
        \draw[fill=k8slightblue, line width=3pt, draw=white] (0,0) circle (2.2);
        % Círculo interior oscuro
        \draw[fill=k8scodebg, line width=2pt, draw=k8scyan] (0,0) circle (1.8);
        % Símbolo de timón de Kubernetes estilizado
        \foreach \angle in {0, 60, 120, 180, 240, 300} {
            \draw[white, line width=4pt, line cap=round] 
                (\angle:0.4) -- (\angle:1.4);
            \draw[k8scyan, line width=2pt, line cap=round] 
                (\angle:0.6) -- (\angle:1.2);
        }
        % Centro del timón
        \draw[fill=white] (0,0) circle (0.3);
        \draw[fill=k8slightblue] (0,0) circle (0.2);
    \end{tikzpicture}
    
    \vspace{0.8cm}
    
    % Título principal con gradiente simulado
    {\fontsize{48}{58}\selectfont\color{white}\textbf{KUBERNETES}}\\[0.3cm]
    {\fontsize{24}{30}\selectfont\color{k8slightblue}\textbf{Certified Kubernetes Administrator}}\\[0.2cm]
    {\fontsize{16}{20}\selectfont\color{k8scyan}\textbf{CKA 2024 | Guía Completa de Estudio}}\\[1.2cm]
    
    % Línea decorativa
    \textcolor{k8slightblue}{\rule{0.7\textwidth}{2pt}}\\[0.8cm]
    
    {\fontsize{18}{22}\selectfont\color{k8slightblue}\textbf{Guía Definitiva de Administración}}\\[0.2cm]
    {\fontsize{14}{17}\selectfont\color{white}\textit{De Principiante a Administrador Certificado}}\\[1.5cm]
    
    % Caja de contenido mejorada
    \begin{center}
    \begin{tcolorbox}[
        colback=k8scodebg!90,
        colframe=k8slightblue,
        boxrule=2pt,
        arc=10pt,
        width=0.85\textwidth,
        shadow={3pt}{3pt}{0pt}{black!50},
        borderline west={5pt}{0pt}{k8scyan}
    ]
        \centering
        {\fontsize{14}{17}\selectfont\color{k8slightblue}\textbf{\faBook\ CONTENIDO INCLUIDO}}\\[0.4cm]
        
        \begin{minipage}{0.45\textwidth}
            \color{white}
            {\small
            \textcolor{k8scyan}{\faServer} \textbf{Arquitectura \& Instalación}\\[0.2cm]
            \textcolor{k8scyan}{\faDatabase} \textbf{Storage \& Persistence}\\[0.2cm]
            \textcolor{k8scyan}{\faNetworkWired} \textbf{Networking \& Services}\\[0.2cm]
            }
        \end{minipage}
        \hfill
        \begin{minipage}{0.45\textwidth}
            \color{white}
            {\small
            \textcolor{k8scyan}{\faCog} \textbf{Workloads \& Scheduling}\\[0.2cm]
            \textcolor{k8scyan}{\faShield} \textbf{Security \& RBAC}\\[0.2cm]
            \textcolor{k8scyan}{\faTools} \textbf{Troubleshooting \& Debug}\\[0.2cm]
            }
        \end{minipage}\\[0.3cm]
        
        \textcolor{k8slightblue}{\rule{0.8\textwidth}{1pt}}\\[0.2cm]
        {\small\color{k8slightblue}\textbf{+ 200 Comandos Prácticos | Ejercicios CKA | Laboratorios}}
    \end{tcolorbox}
    \end{center}
    
    \vspace{1cm}
    
    % Información adicional estilizada
    \begin{tikzpicture}
        \node[
            rectangle,
            rounded corners=5pt,
            fill=k8slightblue!20,
            draw=k8slightblue,
            line width=1pt,
            minimum width=6cm,
            minimum height=1cm
        ] at (0,0) {
            \textcolor{k8snavyblue}{\textbf{\faCheckCircle\ Actualizado 2024 | \faCheckCircle\ Ejemplos Reales | \faCheckCircle\ Best Practices}}
        };
    \end{tikzpicture}
    
    \vfill
    
    {\fontsize{12}{15}\selectfont\color{k8sdarkgray}\textbf{Enero 2026 | Versión 1.30+}}
    
    \restoregeometry
    \nopagecolor
\end{titlepage}

% ====== ÍNDICE ======
\tableofcontents
\newpage

% ====== INTRODUCCIÓN ======
\section{Introducción a Kubernetes}

\begin{infobox}
Kubernetes es una plataforma de orquestación de contenedores de código abierto que automatiza el despliegue, escalado y gestión de aplicaciones en contenedores.
\end{infobox}

\subsection{¿Qué es Kubernetes?}

Kubernetes (también conocido como K8s) es un sistema de orquestación de contenedores que fue desarrollado originalmente por Google y ahora es mantenido por la Cloud Native Computing Foundation (CNCF). Su nombre proviene del griego y significa "timonel" o "piloto".

\subsubsection{Características Principales}

\begin{itemize}[leftmargin=2cm]
    \item[\textcolor{k8sblue}{\faDocker}] \textbf{Orquestación de Contenedores}: Gestiona automáticamente la vida útil de los contenedores
    \item[\textcolor{k8ssuccess}{\faArrowsAltH}] \textbf{Escalado Automático}: Escala aplicaciones hacia arriba o abajo según la demanda
    \item[\textcolor{k8swarning}{\faHeart}] \textbf{Auto-sanación}: Reinicia contenedores fallidos, reemplaza nodos, etc.
    \item[\textcolor{k8ssteelblue}{\faBalanceScale}] \textbf{Load Balancing}: Distribuye el tráfico de red automáticamente
    \item[\textcolor{k8serror}{\faLock}] \textbf{Gestión de Secrets}: Manejo seguro de información sensible
\end{itemize}

\subsection{Arquitectura de Kubernetes}

Kubernetes sigue una arquitectura maestro-esclavo (master-worker) con los siguientes componentes principales:

\subsubsection{Control Plane (Plano de Control)}

\begin{tipbox}
El Control Plane es el cerebro de Kubernetes, donde se toman todas las decisiones de gestión del cluster.
\end{tipbox}

\begin{enumerate}
    \item \textbf{\textcolor{k8sblue}{API Server (kube-apiserver)}}
    \begin{itemize}
        \item Punto de entrada principal para todas las operaciones REST
        \item Valida y configura datos para objetos API
        \item Sirve como frontend para el estado compartido del cluster
    \end{itemize}
    
    \item \textbf{\textcolor{k8ssuccess}{etcd}}
    \begin{itemize}
        \item Base de datos distribuida y consistente
        \item Almacena toda la configuración del cluster
        \item Utiliza el algoritmo de consenso Raft
    \end{itemize}
    
    \item \textbf{\textcolor{k8swarning}{Controller Manager (kube-controller-manager)}}
    \begin{itemize}
        \item Ejecuta los controladores del cluster
        \item Node Controller, Replication Controller, Endpoints Controller
        \item Service Account \& Token Controllers
    \end{itemize}
    
    \item \textbf{\textcolor{k8ssteelblue}{Scheduler (kube-scheduler)}}
    \begin{itemize}
        \item Asigna pods a nodos específicos
        \item Considera recursos, políticas y restricciones
        \item Optimiza la distribución de cargas de trabajo
    \end{itemize}
\end{enumerate}

\subsubsection{Worker Nodes (Nodos de Trabajo)}

\begin{enumerate}
    \item \textbf{\textcolor{k8sblue}{Kubelet}}
    \begin{itemize}
        \item Agente que se ejecuta en cada nodo
        \item Se comunica con el API Server
        \item Gestiona pods y contenedores en el nodo
    \end{itemize}
    
    \item \textbf{\textcolor{k8ssuccess}{Kube-proxy}}
    \begin{itemize}
        \item Componente de red que mantiene reglas de red
        \item Permite comunicación de red a los pods
        \item Implementa parte del concepto de Service
    \end{itemize}
    
    \item \textbf{\textcolor{k8swarning}{Container Runtime}}
    \begin{itemize}
        \item Docker, containerd, CRI-O
        \item Responsable de ejecutar contenedores
        \item Gestiona imágenes de contenedores
    \end{itemize}
\end{enumerate}

\newpage

% ====== ARQUITECTURA, INSTALACIÓN Y CONFIGURACIÓN ======
\section{Cluster Architecture, Installation \& Configuration}

\begin{warningbox}
Esta sección cubre el 25\% del examen CKA y es fundamental para entender cómo funciona Kubernetes internamente.
\end{warningbox}

\subsection{Role Based Access Control (RBAC)}

RBAC es un método para regular el acceso a recursos informáticos basado en los roles de usuarios individuales dentro de una organización.

\subsubsection{Componentes de RBAC}

\begin{enumerate}
    \item \textbf{\textcolor{k8sblue}{Role y ClusterRole}}
    \begin{itemize}
        \item \textbf{Role}: Define permisos dentro de un namespace específico
        \item \textbf{ClusterRole}: Define permisos a nivel de cluster
    \end{itemize}
    
    \item \textbf{\textcolor{k8ssuccess}{RoleBinding y ClusterRoleBinding}}
    \begin{itemize}
        \item \textbf{RoleBinding}: Vincula un Role a usuarios/grupos en un namespace
        \item \textbf{ClusterRoleBinding}: Vincula un ClusterRole a nivel de cluster
    \end{itemize}
    
    \item \textbf{\textcolor{k8swarning}{Subjects}}
    \begin{itemize}
        \item Users (usuarios)
        \item Groups (grupos)
        \item ServiceAccounts (cuentas de servicio)
    \end{itemize}
\end{enumerate}

\subsubsection{Ejemplo Práctico de RBAC}

\begin{commandbox}
\# Crear un Role para leer pods
apiVersion: rbac.authorization.k8s.io/v1
kind: Role
metadata:
  namespace: default
  name: pod-reader
rules:
- apiGroups: [""]
  resources: ["pods"]
  verbs: ["get", "watch", "list"]
\end{commandbox}

\begin{commandbox}
\# Crear un RoleBinding
apiVersion: rbac.authorization.k8s.io/v1
kind: RoleBinding
metadata:
  name: read-pods
  namespace: default
subjects:
- kind: User
  name: jane
  apiGroup: rbac.authorization.k8s.io
roleRef:
  kind: Role
  name: pod-reader
  apiGroup: rbac.authorization.k8s.io
\end{commandbox}

\subsubsection{Comandos Útiles para RBAC}

\begin{commandbox}
\# Verificar permisos de un usuario
kubectl auth can-i get pods --as=jane

\# Verificar todos los permisos
kubectl auth can-i --list --as=jane

\# Crear ServiceAccount
kubectl create serviceaccount my-service-account

\# Obtener token de ServiceAccount
kubectl create token my-service-account
\end{commandbox}

\subsection{Instalación de Kubernetes con kubeadm}

kubeadm es una herramienta oficial para crear clusters de Kubernetes que funciona rápida y fácilmente.

\subsubsection{Preparación del Sistema}

\begin{commandbox}
\# Deshabilitar swap
sudo swapoff -a
sudo sed -i 's|/ swap |/# swap |g' /etc/fstab

\# Cargar módulos del kernel necesarios
cat <<EOF | sudo tee /etc/modules-load.d/k8s.conf
overlay
br\_netfilter
EOF

sudo modprobe overlay
sudo modprobe br\_netfilter

\# Configurar parámetros del kernel
cat <<EOF | sudo tee /etc/sysctl.d/k8s.conf
net.bridge.bridge-nf-call-iptables  = 1
net.bridge.bridge-nf-call-ip6tables = 1
net.ipv4.ip\_forward                 = 1
EOF

sudo sysctl --system
\end{commandbox}

\subsubsection{Instalación de Container Runtime (containerd)}

\begin{commandbox}
\# Instalar containerd
sudo apt-get update
sudo apt-get install -y containerd

\# Configurar containerd
sudo mkdir -p /etc/containerd
containerd config default | sudo tee /etc/containerd/config.toml

\# Editar config.toml para usar systemd como cgroup driver
sudo sed -i 's/SystemdCgroup = false/SystemdCgroup = true/' \\
  /etc/containerd/config.toml

\# Reiniciar containerd
sudo systemctl restart containerd
sudo systemctl enable containerd
\end{commandbox}

\subsubsection{Instalación de kubeadm, kubelet y kubectl}

\begin{commandbox}
\# Agregar repositorio de Kubernetes
sudo apt-get update
sudo apt-get install -y apt-transport-https ca-certificates curl

curl -fsSL https://pkgs.k8s.io/core:/stable:/v1.29/deb/Release.key | \\
  sudo gpg --dearmor -o /etc/apt/keyrings/kubernetes-apt-keyring.gpg

echo 'deb [signed-by=/etc/apt/keyrings/kubernetes-apt-keyring.gpg] \\
  https://pkgs.k8s.io/core:/stable:/v1.29/deb/ /' | \\
  sudo tee /etc/apt/sources.list.d/kubernetes.list

\# Instalar componentes
sudo apt-get update
sudo apt-get install -y kubelet kubeadm kubectl
sudo apt-mark hold kubelet kubeadm kubectl

\# Habilitar kubelet
sudo systemctl enable kubelet
\end{commandbox}

\subsubsection{Inicialización del Cluster}

\begin{commandbox}
\# Inicializar el nodo maestro
sudo kubeadm init --pod-network-cidr=10.244.0.0/16

\# Configurar kubectl para el usuario
mkdir -p \$HOME/.kube
sudo cp -i /etc/kubernetes/admin.conf \$HOME/.kube/config
sudo chown $(id -u):$(id -g) \$HOME/.kube/config

\# Instalar plugin de red (Flannel)
kubectl apply -f \\
  https://github.com/flannel-io/flannel/releases/latest/download/kube-flannel.yml
\end{commandbox}

\begin{tipbox}
Guarda el comando \texttt{kubeadm join} que aparece al final de la inicialización. Lo necesitarás para unir nodos worker al cluster.
\end{tipbox}

\subsection{Alta Disponibilidad en Kubernetes}

Un cluster de alta disponibilidad garantiza que el plano de control esté replicado a través de múltiples nodos para evitar puntos únicos de falla.

\subsubsection{Topologías de Alta Disponibilidad}

\begin{enumerate}
    \item \textbf{\textcolor{k8sblue}{Stacked etcd}}
    \begin{itemize}
        \item etcd se ejecuta en los mismos nodos que el control plane
        \item Menos infraestructura requerida
        \item Mayor riesgo de pérdida de datos
    \end{itemize}
    
    \item \textbf{\textcolor{k8ssuccess}{External etcd}}
    \begin{itemize}
        \item etcd se ejecuta en nodos separados
        \item Mayor resiliencia
        \item Más infraestructura requerida
    \end{itemize}
\end{enumerate}

\subsubsection{Configuración con Load Balancer}

\begin{commandbox}
\# Configuración de HAProxy para load balancer
\# /etc/haproxy/haproxy.cfg

global
    maxconn 4096
    log stdout local0

defaults
    mode tcp
    timeout connect 5000ms
    timeout client 50000ms
    timeout server 50000ms

frontend k8s-api
    bind *:6443
    default\_backend k8s-masters

backend k8s-masters
    balance roundrobin
    server master1 192.168.1.10:6443 check
    server master2 192.168.1.11:6443 check
    server master3 192.168.1.12:6443 check
\end{commandbox}

\subsubsection{Inicialización de HA Cluster}

\begin{commandbox}
\# En el primer nodo maestro
sudo kubeadm init --control-plane-endpoint="lb.example.com:6443" \\
  --upload-certs --pod-network-cidr=10.244.0.0/16

\# Para unir nodos maestros adicionales
sudo kubeadm join lb.example.com:6443 --token <token> \\
  --discovery-token-ca-cert-hash sha256:<hash> \\
  --control-plane --certificate-key <key>

\# Para unir nodos worker
sudo kubeadm join lb.example.com:6443 --token <token> \\
  --discovery-token-ca-cert-hash sha256:<hash>
\end{commandbox}

\subsection{Actualización de Clusters con kubeadm}

\subsubsection{Planificación de la Actualización}

\begin{warningbox}
Siempre realizar copias de seguridad de etcd antes de actualizar. Las actualizaciones deben ser secuenciales (v1.28 → v1.29 → v1.30).
\end{warningbox}

\begin{commandbox}
\# Verificar versión actual
kubectl get nodes

\# Verificar versiones disponibles
apt list -a kubeadm

\# Planificar actualización
sudo kubeadm upgrade plan
\end{commandbox}

\subsubsection{Actualización del Control Plane}

\begin{commandbox}
\# Actualizar kubeadm
sudo apt-mark unhold kubeadm
sudo apt-get update \ && \ &&  sudo apt-get install -y kubeadm=1.29.0-00
sudo apt-mark hold kubeadm

\# Verificar la nueva versión
kubeadm version

\# Realizar la actualización
sudo kubeadm upgrade apply v1.29.0

\# Drenar el nodo
kubectl drain <node-name> --ignore-daemonsets

\# Actualizar kubelet y kubectl
sudo apt-mark unhold kubelet kubectl
sudo apt-get update \ && \ &&  sudo apt-get install -y \\
  kubelet=1.29.0-00 kubectl=1.29.0-00
sudo apt-mark hold kubelet kubectl

\# Reiniciar kubelet
sudo systemctl daemon-reload
sudo systemctl restart kubelet

\# Quitar el taint del nodo
kubectl uncordon <node-name>
\end{commandbox}

\subsubsection{Actualización de Worker Nodes}

\begin{commandbox}
\# En el nodo worker
\# Actualizar kubeadm
sudo apt-mark unhold kubeadm
sudo apt-get update \ && \ &&  sudo apt-get install -y kubeadm=1.29.0-00
sudo apt-mark hold kubeadm

\# Actualizar configuración del nodo
sudo kubeadm upgrade node

\# Drenar el nodo (desde el control plane)
kubectl drain <worker-node-name> --ignore-daemonsets \\
  --delete-emptydir-data

\# Actualizar kubelet y kubectl
sudo apt-mark unhold kubelet kubectl
sudo apt-get update \ && \ &&  sudo apt-get install -y \\
  kubelet=1.29.0-00 kubectl=1.29.0-00
sudo apt-mark hold kubelet kubectl

\# Reiniciar kubelet
sudo systemctl daemon-reload
sudo systemctl restart kubelet

\# Quitar el taint del nodo (desde el control plane)
kubectl uncordon <worker-node-name>
\end{commandbox}

\subsection{Backup y Restore de etcd}

etcd es el almacén de datos del cluster, por lo que es crucial mantener copias de seguridad regulares.

\subsubsection{Backup de etcd}

\begin{commandbox}
\# Verificar el endpoint de etcd
kubectl get pod etcd-master -n kube-system -o yaml | grep -A1 listen-client

\# Realizar backup
ETCDCTL\_API=3 etcdctl snapshot save /tmp/etcd-backup.db \\
  --endpoints=https://127.0.0.1:2379 \\
  --cacert=/etc/kubernetes/pki/etcd/ca.crt \\
  --cert=/etc/kubernetes/pki/etcd/server.crt \\
  --key=/etc/kubernetes/pki/etcd/server.key

\# Verificar el backup
ETCDCTL\_API=3 etcdctl snapshot status /tmp/etcd-backup.db \\
  --write-out=table
\end{commandbox}

\subsubsection{Restore de etcd}

\begin{errorbox}
El proceso de restore requiere detener el control plane y puede causar tiempo de inactividad.
\end{errorbox}

\begin{commandbox}
\# Detener API server y etcd
sudo mv /etc/kubernetes/manifests /etc/kubernetes/manifests.backup

\# Restaurar desde backup
ETCDCTL\_API=3 etcdctl snapshot restore /tmp/etcd-backup.db \\
  --data-dir=/var/lib/etcd-restore \\
  --initial-cluster=master=https://192.168.1.10:2380 \\
  --initial-advertise-peer-urls=https://192.168.1.10:2380

\# Actualizar el manifest de etcd
sudo sed -i 's|/var/lib/etcd|/var/lib/etcd-restore|' \\
  /etc/kubernetes/manifests.backup/etcd.yaml

\# Restaurar manifests
sudo mv /etc/kubernetes/manifests.backup /etc/kubernetes/manifests

\# Verificar que el cluster está funcionando
kubectl get nodes
\end{commandbox}

\newpage

% ====== STORAGE, SERVICES & NETWORKING ======
\section{Storage, Services \& Networking}

\begin{infobox}
Esta sección representa el 20\% del examen CKA y cubre aspectos críticos de almacenamiento y networking en Kubernetes.
\end{infobox}

\subsection{Storage en Kubernetes}

El almacenamiento en Kubernetes permite que las aplicaciones persistan datos más allá del ciclo de vida de los pods.

\subsubsection{Tipos de Volúmenes}

\begin{enumerate}
    \item \textbf{\textcolor{k8sblue}{EmptyDir}}
    \begin{itemize}
        \item Volumen temporal que existe mientras el pod existe
        \item Se crea cuando un pod se asigna a un nodo
        \item Se elimina cuando el pod se remueve del nodo
    \end{itemize}
    
    \item \textbf{\textcolor{k8ssuccess}{HostPath}}
    \begin{itemize}
        \item Monta un archivo o directorio del nodo host
        \item Útil para acceso a logs del sistema o socket de Docker
        \item No recomendado para aplicaciones en producción
    \end{itemize}
    
    \item \textbf{\textcolor{k8swarning}{PersistentVolume (PV)}}
    \begin{itemize}
        \item Recurso de almacenamiento en el cluster
        \item Independiente del ciclo de vida del pod
        \item Provisionado por un administrador o dinámicamente
    \end{itemize}
\end{enumerate}

\subsubsection{Ejemplo de EmptyDir}

\begin{commandbox}
apiVersion: v1
kind: Pod
metadata:
  name: test-pod
spec:
  containers:
  - name: test-container
    image: nginx
    volumeMounts:
    - name: cache-volume
      mountPath: /cache
  volumes:
  - name: cache-volume
    emptyDir: {}
\end{commandbox}

\subsubsection{Persistent Volumes (PV)}

Los Persistent Volumes son recursos de almacenamiento en el cluster que han sido provisionados por un administrador.

\textbf{Ciclo de Vida de un PV:}

\begin{enumerate}
    \item \textbf{\textcolor{k8sblue}{Provisioning}}: Creación estática o dinámica
    \item \textbf{\textcolor{k8ssuccess}{Binding}}: PVC se vincula a un PV
    \item \textbf{\textcolor{k8swarning}{Using}}: Pod usa el PVC
    \item \textbf{\textcolor{k8serror}{Reclaiming}}: Liberación del PV
\end{enumerate}

\begin{commandbox}
apiVersion: v1
kind: PersistentVolume
metadata:
  name: pv-example
spec:
  capacity:
    storage: 10Gi
  accessModes:
    - ReadWriteOnce
  persistentVolumeReclaimPolicy: Retain
  storageClassName: manual
  hostPath:
    path: "/mnt/data"
\end{commandbox}

\subsubsection{Persistent Volume Claims (PVC)}

Un PVC es una solicitud de almacenamiento por parte de un usuario.

\textbf{Modos de Acceso:}
\begin{itemize}
    \item \textbf{ReadWriteOnce (RWO)}: Montaje de lectura-escritura por un solo nodo
    \item \textbf{ReadOnlyMany (ROX)}: Montaje de solo lectura por múltiples nodos
    \item \textbf{ReadWriteMany (RWX)}: Montaje de lectura-escritura por múltiples nodos
\end{itemize}

\begin{commandbox}
apiVersion: v1
kind: PersistentVolumeClaim
metadata:
  name: pvc-example
spec:
  accessModes:
    - ReadWriteOnce
  resources:
    requests:
      storage: 5Gi
  storageClassName: manual
\end{commandbox}

\subsubsection{Storage Classes}

Las Storage Classes proporcionan una forma de describir las "clases" de almacenamiento disponibles.

\begin{commandbox}
apiVersion: storage.k8s.io/v1
kind: StorageClass
metadata:
  name: fast-ssd
provisioner: kubernetes.io/aws-ebs
parameters:
  type: gp3
  iops: "3000"
  throughput: "125"
reclaimPolicy: Delete
allowVolumeExpansion: true
mountOptions:
  - debug
volumeBindingMode: WaitForFirstConsumer
\end{commandbox}

\begin{tipbox}
El \texttt{volumeBindingMode: WaitForFirstConsumer} es útil para topologías donde el PV debe crearse en la misma zona que el pod.
\end{tipbox}

\subsection{Services en Kubernetes}

Los Services exponen aplicaciones que se ejecutan en un conjunto de pods como un servicio de red.

\subsubsection{Tipos de Services}

\begin{enumerate}
    \item \textbf{\textcolor{k8sblue}{ClusterIP}}
    \begin{itemize}
        \item Expone el servicio en una IP interna del cluster
        \item Solo accesible desde dentro del cluster
        \item Es el tipo por defecto
    \end{itemize}
    
    \item \textbf{\textcolor{k8ssuccess}{NodePort}}
    \begin{itemize}
        \item Expone el servicio en un puerto estático de cada nodo
        \item Accesible desde fuera del cluster usando <NodeIP>:<NodePort>
        \item Puerto en rango 30000-32767
    \end{itemize}
    
    \item \textbf{\textcolor{k8swarning}{LoadBalancer}}
    \begin{itemize}
        \item Expone el servicio usando un load balancer externo
        \item Requiere proveedor de nube
        \item Asigna una IP externa fija
    \end{itemize}
    
    \item \textbf{\textcolor{k8ssteelblue}{ExternalName}}
    \begin{itemize}
        \item Mapea el servicio a un nombre DNS externo
        \item No usa proxying, retorna un registro CNAME
    \end{itemize}
\end{enumerate}

\subsubsection{Ejemplo de Service ClusterIP}

\begin{commandbox}
apiVersion: v1
kind: Service
metadata:
  name: my-service
spec:
  selector:
    app: my-app
  ports:
    - protocol: TCP
      port: 80
      targetPort: 8080
  type: ClusterIP
\end{commandbox}

\subsubsection{Ejemplo de Service NodePort}

\begin{commandbox}
apiVersion: v1
kind: Service
metadata:
  name: my-nodeport-service
spec:
  selector:
    app: my-app
  ports:
    - protocol: TCP
      port: 80
      targetPort: 8080
      nodePort: 30080
  type: NodePort
\end{commandbox}

\subsubsection{Endpoints}

Los Endpoints definen las IPs y puertos de los pods que respaldan un servicio.

\begin{commandbox}
\# Ver endpoints de un servicio
kubectl get endpoints my-service

\# Describir endpoints
kubectl describe endpoints my-service

\# Endpoints manual (sin selector)
apiVersion: v1
kind: Endpoints
metadata:
  name: my-service
subsets:
  - addresses:
      - ip: 192.168.1.100
    ports:
      - port: 8080
\end{commandbox}

\subsection{Ingress Controllers e Ingress Resources}

Ingress gestiona el acceso externo a los servicios en un cluster, típicamente HTTP.

\subsubsection{Conceptos de Ingress}

\begin{itemize}
    \item \textbf{Ingress Controller}: Implementación que satisface las reglas de Ingress
    \item \textbf{Ingress Resource}: Configuración que define las reglas de routing
\end{itemize}

\subsubsection{Instalación de NGINX Ingress Controller}

\begin{commandbox}
\# Instalar NGINX Ingress Controller
kubectl apply -f \\
  https://raw.githubusercontent.com/kubernetes/ingress-nginx/controller-v1.8.2/deploy/static/provider/cloud/deploy.yaml

\# Verificar instalación
kubectl get pods -n ingress-nginx

\# Ver servicio del ingress controller
kubectl get svc -n ingress-nginx
\end{commandbox}

\subsubsection{Ejemplo de Ingress Resource}

\begin{commandbox}
apiVersion: networking.k8s.io/v1
kind: Ingress
metadata:
  name: example-ingress
  annotations:
    nginx.ingress.kubernetes.io/rewrite-target: /
spec:
  ingressClassName: nginx
  rules:
  - host: example.com
    http:
      paths:
      - path: /app1
        pathType: Prefix
        backend:
          service:
            name: app1-service
            port:
              number: 80
      - path: /app2
        pathType: Prefix
        backend:
          service:
            name: app2-service
            port:
              number: 80
\end{commandbox}

\subsubsection{Ingress con TLS}

\begin{commandbox}
apiVersion: networking.k8s.io/v1
kind: Ingress
metadata:
  name: tls-ingress
spec:
  ingressClassName: nginx
  tls:
  - hosts:
    - secure.example.com
    secretName: tls-secret
  rules:
  - host: secure.example.com
    http:
      paths:
      - path: /
        pathType: Prefix
        backend:
          service:
            name: secure-service
            port:
              number: 443
\end{commandbox}

\subsection{CoreDNS}

CoreDNS es el servidor DNS del cluster que proporciona resolución de nombres para servicios.

\subsubsection{Configuración de CoreDNS}

\begin{commandbox}
\# Ver configuración de CoreDNS
kubectl get configmap coredns -n kube-system -o yaml

\# Configuración típica
.:53 {
    errors
    health {
       lameduck 5s
    }
    ready
    kubernetes cluster.local in-addr.arpa ip6.arpa {
       pods insecure
       fallthrough in-addr.arpa ip6.arpa
       ttl 30
    }
    prometheus :9153
    forward . /etc/resolv.conf
    cache 30
    loop
    reload
    loadbalance
}
\end{commandbox}

\subsubsection{Resolución DNS en Kubernetes}

\begin{itemize}
    \item \textbf{Services}: \texttt{my-service.my-namespace.svc.cluster.local}
    \item \textbf{Pods}: \texttt{192-168-1-1.my-namespace.pod.cluster.local}
    \item \textbf{Dentro del mismo namespace}: \texttt{my-service}
\end{itemize}

\begin{commandbox}
\# Probar resolución DNS desde un pod
kubectl run test-dns --image=busybox:1.28 --rm -it --restart=Never \\
  -- nslookup kubernetes.default.svc.cluster.local

\# Verificar configuración DNS de un pod
kubectl exec test-pod -- cat /etc/resolv.conf
\end{commandbox}

\subsubsection{Troubleshooting DNS}

\begin{errorbox}
Problemas comunes: CoreDNS pods no running, configuración incorrecta de iptables, problemas de conectividad de red.
\end{errorbox}

\begin{commandbox}
\# Verificar estado de CoreDNS
kubectl get pods -n kube-system -l k8s-app=kube-dns

\# Ver logs de CoreDNS
kubectl logs -n kube-system deployment/coredns

\# Probar conectividad DNS
kubectl run dns-test --image=busybox:1.28 --rm -it --restart=Never \\
  -- sh -c 'nslookup kubernetes.default \ && \ &&  echo "DNS OK"'
\end{commandbox}

\subsection{Container Network Interface (CNI)}

CNI define plugins para configurar interfaces de red en contenedores Linux.

\subsubsection{Plugins CNI Populares}

\begin{enumerate}
    \item \textbf{\textcolor{k8sblue}{Flannel}}
    \begin{itemize}
        \item Simple y fácil de configurar
        \item Usa VXLAN por defecto
        \item Bueno para comenzar
    \end{itemize}
    
    \item \textbf{\textcolor{k8ssuccess}{Calico}}
    \begin{itemize}
        \item Soporta políticas de red
        \item Usa BGP para routing
        \item Alta performance
    \end{itemize}
    
    \item \textbf{\textcolor{k8swarning}{Weave}}
    \begin{itemize}
        \item Mesh network automático
        \item Encryption built-in
        \item Easy troubleshooting
    \end{itemize}
\end{enumerate}

\subsubsection{Instalación de Calico}

\begin{commandbox}
\# Instalar Calico
kubectl create -f \\
  https://raw.githubusercontent.com/projectcalico/calico/v3.26.1/manifests/tigera-operator.yaml

kubectl create -f \\
  https://raw.githubusercontent.com/projectcalico/calico/v3.26.1/manifests/custom-resources.yaml

\# Verificar instalación
kubectl get pods -n calico-system

\# Ver nodos Calico
kubectl get nodes -o wide
\end{commandbox}

\subsubsection{Network Policies con Calico}

\begin{commandbox}
\# Denegar todo el tráfico por defecto
apiVersion: networking.k8s.io/v1
kind: NetworkPolicy
metadata:
  name: default-deny-all
  namespace: production
spec:
  podSelector: {}
  policyTypes:
  - Ingress
  - Egress

\# Permitir tráfico específico
apiVersion: networking.k8s.io/v1
kind: NetworkPolicy
metadata:
  name: allow-frontend-to-backend
spec:
  podSelector:
    matchLabels:
      role: backend
  policyTypes:
  - Ingress
  ingress:
  - from:
    - podSelector:
        matchLabels:
          role: frontend
    ports:
    - protocol: TCP
      port: 8080
\end{commandbox}

\newpage

% ====== WORKLOADS & SCHEDULING ======
\section{Workloads \& Scheduling}

\begin{warningbox}
Esta sección constituye el 20\% del examen CKA y cubre la gestión de aplicaciones y programación de pods.
\end{warningbox}

\subsection{Pods}

Los Pods son la unidad básica de despliegue en Kubernetes, conteniendo uno o más contenedores que comparten almacenamiento y red.

\subsubsection{Anatomía de un Pod}

\begin{commandbox}
apiVersion: v1
kind: Pod
metadata:
  name: my-pod
  labels:
    app: my-app
    version: v1.0
  annotations:
    description: "Mi aplicación principal"
spec:
  containers:
  - name: main-container
    image: nginx:1.20
    ports:
    - containerPort: 80
      name: http
    env:
    - name: ENV_VAR
      value: "production"
    resources:
      requests:
        memory: "64Mi"
        cpu: "250m"
      limits:
        memory: "128Mi"
        cpu: "500m"
    volumeMounts:
    - name: config-volume
      mountPath: /etc/nginx/conf.d
    livenessProbe:
      httpGet:
        path: /health
        port: 80
      initialDelaySeconds: 30
      periodSeconds: 10
    readinessProbe:
      httpGet:
        path: /ready
        port: 80
      initialDelaySeconds: 5
      periodSeconds: 5
  volumes:
  - name: config-volume
    configMap:
      name: nginx-config
  restartPolicy: Always
  nodeSelector:
    disktype: ssd
\end{commandbox}

\subsubsection{Ciclo de Vida de un Pod}

\begin{enumerate}
    \item \textbf{\textcolor{k8sblue}{Pending}}: Pod ha sido aceptado pero no todos los contenedores están running
    \item \textbf{\textcolor{k8ssuccess}{Running}}: Pod está ejecutándose en un nodo
    \item \textbf{\textcolor{k8swarning}{Succeeded}}: Todos los contenedores han terminado exitosamente
    \item \textbf{\textcolor{k8serror}{Failed}}: Todos los contenedores han terminado, al menos uno falló
    \item \textbf{\textcolor{k8ssteelblue}{Unknown}}: Estado del pod no puede ser determinado
\end{enumerate}

\subsubsection{Comandos Útiles para Pods}

\begin{commandbox}
\# Crear un pod imperativo
kubectl run nginx --image=nginx:1.20 --port=80

\# Obtener pods con más información
kubectl get pods -o wide

\# Describir un pod
kubectl describe pod my-pod

\# Ver logs de un pod
kubectl logs my-pod -f

\# Ejecutar comando en un pod
kubectl exec my-pod -it -- /bin/bash

\# Reenviar puerto
kubectl port-forward pod/my-pod 8080:80

\# Eliminar un pod
kubectl delete pod my-pod
\end{commandbox}

\subsection{Deployments}

Los Deployments proporcionan actualizaciones declarativas para Pods y ReplicaSets.

\subsubsection{Ejemplo de Deployment}

\begin{commandbox}
apiVersion: apps/v1
kind: Deployment
metadata:
  name: nginx-deployment
  labels:
    app: nginx
spec:
  replicas: 3
  selector:
    matchLabels:
      app: nginx
  template:
    metadata:
      labels:
        app: nginx
    spec:
      containers:
      - name: nginx
        image: nginx:1.20
        ports:
        - containerPort: 80
        resources:
          requests:
            cpu: 100m
            memory: 128Mi
          limits:
            cpu: 200m
            memory: 256Mi
  strategy:
    type: RollingUpdate
    rollingUpdate:
      maxSurge: 1
      maxUnavailable: 1
\end{commandbox}
\end{tcolorbox}
\subsubsection{Rolling Updates y Rollbacks}

\begin{commandbox}
\# Actualizar imagen del deployment
kubectl set image deployment/nginx-deployment nginx=nginx:1.21

\# Ver estado del rollout
kubectl rollout status deployment/nginx-deployment

\# Ver historial de rollouts
kubectl rollout history deployment/nginx-deployment

\# Hacer rollback a versión anterior
kubectl rollout undo deployment/nginx-deployment

\# Hacer rollback a versión específica
kubectl rollout undo deployment/nginx-deployment --to-revision=2

\# Pausar un rollout
kubectl rollout pause deployment/nginx-deployment

\# Resumir un rollout
kubectl rollout resume deployment/nginx-deployment
\end{commandbox}

\subsubsection{Estrategias de Deployment}

\begin{enumerate}
    \item \textbf{\textcolor{k8sblue}{RollingUpdate}}
    \begin{itemize}
        \item Reemplaza pods gradualmente
        \item Zero downtime
        \item Configurable con maxSurge y maxUnavailable
    \end{itemize}
    
    \item \textbf{\textcolor{k8ssuccess}{Recreate}}
    \begin{itemize}
        \item Elimina todos los pods antes de crear nuevos
        \item Causa downtime
        \item Útil cuando no se pueden ejecutar versiones múltiples
    \end{itemize}
\end{enumerate}

\subsection{ConfigMaps y Secrets}

ConfigMaps y Secrets permiten separar la configuración del código de la aplicación.

\subsubsection{ConfigMaps}

Los ConfigMaps almacenan datos de configuración no confidenciales en pares clave-valor.

\begin{commandbox}
\# Crear ConfigMap desde línea de comandos
kubectl create configmap app-config \\
  --from-literal=database\_host=localhost \\
  --from-literal=database_port=5432

\# Crear ConfigMap desde archivo
kubectl create configmap nginx-config --from-file=nginx.conf

\# ConfigMap declarativo
apiVersion: v1
kind: ConfigMap
metadata:
  name: game-config
data:
  game.properties: |
    enemies=aliens
    lives=3
    enemies.cheat=true
  ui.properties: |
    color.good=purple
    color.bad=yellow
    allow.textmode=true
\end{commandbox}

\subsubsection{Usar ConfigMaps en Pods}

\begin{commandbox}
apiVersion: v1
kind: Pod
metadata:
  name: configmap-pod
spec:
  containers:
  - name: test-container
    image: busybox
    command: [ "/bin/sh", "-c", "env" ]
    env:
    # Variable desde ConfigMap
    - name: SPECIAL\_LEVEL\_KEY
      valueFrom:
        configMapKeyRef:
          name: special-config
          key: special.how
    # Todas las claves como variables de entorno
    envFrom:
    - configMapRef:
        name: game-config
    volumeMounts:
    - name: config-volume
      mountPath: /etc/config
  volumes:
  - name: config-volume
    configMap:
      name: game-config
\end{commandbox}

\subsubsection{Secrets}

Los Secrets almacenan información sensible como passwords, tokens OAuth, y claves SSH.

\begin{commandbox}
\# Crear Secret desde línea de comandos
kubectl create secret generic db-secret \\
  --from-literal=username=admin \\
  --from-literal=password=secretpassword

\# Crear Secret para Docker registry
kubectl create secret docker-registry regcred \\
  --docker-server=myregistry.com \\
  --docker-username=myuser \\
  --docker-password=mypassword \\
  --docker-email=myemail@example.com

\# Secret declarativo
apiVersion: v1
kind: Secret
metadata:
  name: mysecret
type: Opaque
data:
  username: YWRtaW4=  # base64 encoded
  password: MWYyZDFlMmU2N2Rm  # base64 encoded
\end{commandbox}

\begin{tipbox}
Los valores en Secrets deben estar codificados en base64. Usa \texttt{echo -n 'mypassword' | base64} para codificar.
\end{tipbox}

\subsubsection{Tipos de Secrets}

\begin{itemize}
    \item \textbf{Opaque}: Datos de usuario arbitrarios
    \item \textbf{kubernetes.io/service-account-token}: Token de cuenta de servicio
    \item \textbf{kubernetes.io/dockercfg}: Archivo .dockercfg serializado
    \item \textbf{kubernetes.io/tls}: Datos para un cliente o servidor TLS
\end{itemize}

\subsection{Scaling de Aplicaciones}

\subsubsection{Scaling Manual}

\begin{commandbox}
\# Escalar deployment
kubectl scale deployment nginx-deployment --replicas=5

\# Escalar replicaset
kubectl scale rs my-replicaset --replicas=3

\# Escalar usando archivo
kubectl scale --replicas=3 -f deployment.yaml
\end{commandbox}

\subsubsection{Horizontal Pod Autoscaler (HPA)}

El HPA escala automáticamente el número de pods basado en CPU, memoria u otras métricas.

\begin{commandbox}
\# Crear HPA basado en CPU
kubectl autoscale deployment nginx-deployment --cpu-percent=50 --min=1 --max=10

\# HPA declarativo
apiVersion: autoscaling/v2
kind: HorizontalPodAutoscaler
metadata:
  name: nginx-hpa
spec:
  scaleTargetRef:
    apiVersion: apps/v1
    kind: Deployment
    name: nginx-deployment
  minReplicas: 2
  maxReplicas: 10
  metrics:
  - type: Resource
    resource:
      name: cpu
      target:
        type: Utilization
        averageUtilization: 70
  - type: Resource
    resource:
      name: memory
      target:
        type: Utilization
        averageUtilization: 80
\end{commandbox}

\begin{warningbox}
Para que HPA funcione, necesitas tener el Metrics Server instalado y los pods deben tener resource requests definidos.
\end{warningbox}

\subsubsection{Vertical Pod Autoscaler (VPA)}

VPA ajusta automáticamente los resource requests y limits de los contenedores.

\begin{commandbox}
apiVersion: autoscaling.k8s.io/v1
kind: VerticalPodAutoscaler
metadata:
  name: nginx-vpa
spec:
  targetRef:
    apiVersion: apps/v1
    kind: Deployment
    name: nginx-deployment
  updatePolicy:
    updateMode: "Auto"
  resourcePolicy:
    containerPolicies:
    - containerName: nginx
      maxAllowed:
        cpu: 1
        memory: 500Mi
      minAllowed:
        cpu: 100m
        memory: 50Mi
\end{commandbox}

\subsection{Pod Scheduling}

El scheduling determina en qué nodos se ejecutan los pods.

\subsubsection{NodeSelector}

La forma más simple de restricrir pods a nodos específicos.

\begin{commandbox}
\# Etiquetar un nodo
kubectl label nodes worker-1 disktype=ssd

\# Pod con nodeSelector
apiVersion: v1
kind: Pod
metadata:
  name: nginx
spec:
  containers:
  - name: nginx
    image: nginx
  nodeSelector:
    disktype: ssd
\end{commandbox}

\subsubsection{Node Affinity}

Proporciona reglas más expresivas para seleccionar nodos.

\begin{commandbox}
apiVersion: v1
kind: Pod
metadata:
  name: with-node-affinity
spec:
  affinity:
    nodeAffinity:
      requiredDuringSchedulingIgnoredDuringExecution:
        nodeSelectorTerms:
        - matchExpressions:
          - key: kubernetes.io/arch
            operator: In
            values:
            - amd64
            - arm64
      preferredDuringSchedulingIgnoredDuringExecution:
      - weight: 1
        preference:
          matchExpressions:
          - key: disktype
            operator: In
            values:
            - ssd
  containers:
  - name: with-node-affinity
    image: nginx
\end{commandbox}

\subsubsection{Pod Affinity y Anti-Affinity}

Permite programar pods basándose en qué otros pods ya están ejecutándose en el nodo.

\begin{commandbox}
apiVersion: apps/v1
kind: Deployment
metadata:
  name: redis-cache
spec:
  selector:
    matchLabels:
      app: store
  replicas: 3
  template:
    metadata:
      labels:
        app: store
    spec:
      affinity:
        podAntiAffinity:
          requiredDuringSchedulingIgnoredDuringExecution:
          - labelSelector:
              matchExpressions:
              - key: app
                operator: In
                values:
                - store
            topologyKey: "kubernetes.io/hostname"
        podAffinity:
          preferredDuringSchedulingIgnoredDuringExecution:
          - weight: 100
            podAffinityTerm:
              labelSelector:
                matchExpressions:
                - key: security
                  operator: In
                  values:
                  - S1
              topologyKey: topology.kubernetes.io/zone
      containers:
      - name: web-app
        image: nginx:1.16-alpine
\end{commandbox}

\subsubsection{Taints y Tolerations}

Los Taints permiten a un nodo repeler un conjunto de pods, mientras que las Tolerations permiten a los pods programarse en nodos con taints correspondientes.

\begin{commandbox}
\# Aplicar taint a un nodo
kubectl taint nodes node1 key1=value1:NoSchedule

\# Remover taint
kubectl taint nodes node1 key1=value1:NoSchedule-

\# Pod con tolerations
apiVersion: v1
kind: Pod
metadata:
  name: nginx
spec:
  containers:
  - name: nginx
    image: nginx
  tolerations:
  - key: "key1"
    operator: "Equal"
    value: "value1"
    effect: "NoSchedule"
  - key: "key2"
    operator: "Exists"
    effect: "NoExecute"
    tolerationSeconds: 3600
\end{commandbox}

\subsubsection{Resource Limits y Requests}

Definen cuántos recursos puede usar un contenedor.

\begin{commandbox}
apiVersion: v1
kind: Pod
metadata:
  name: resource-demo
spec:
  containers:
  - name: resource-demo
    image: nginx
    resources:
      requests:
        memory: "64Mi"
        cpu: "250m"
        ephemeral-storage: "1Gi"
      limits:
        memory: "128Mi"
        cpu: "500m"
        ephemeral-storage: "2Gi"
\end{commandbox}

\begin{tipbox}
\textbf{Requests}: Recursos garantizados\\
\textbf{Limits}: Recursos máximos que puede usar\\
CPU se especifica en millicores (m) o cores\\
Memoria en bytes (Mi, Gi, etc.)
\end{tipbox}

\subsubsection{Quality of Service (QoS)}

Kubernetes asigna una clase QoS a cada pod:

\begin{enumerate}
    \item \textbf{\textcolor{k8ssuccess}{Guaranteed}}: Requests = Limits para todos los contenedores
    \item \textbf{\textcolor{k8swarning}{Burstable}}: Al menos un contenedor tiene requests pero no guaranteed
    \item \textbf{\textcolor{k8serror}{BestEffort}}: Ningún contenedor tiene requests ni limits
\end{enumerate}

\newpage

% Continuar con el resto de las secciones...
% (Security, Troubleshooting, etc.)

\section{Security}

\begin{errorbox}
La seguridad en Kubernetes es crítica. Esta sección cubre el 15\% del examen CKA y incluye autenticación, autorización y políticas de seguridad.
\end{errorbox}

\subsection{Autenticación en Kubernetes}

Kubernetes soporta múltiples métodos de autenticación para usuarios y service accounts.

\subsubsection{Métodos de Autenticación}

\begin{enumerate}
    \item \textbf{\textcolor{k8sblue}{X.509 Certificates}}
    \begin{itemize}
        \item Más común para usuarios administradores
        \item Cliente presenta certificado válido
        \item CN del certificado = username, O = groups
    \end{itemize}
    
    \item \textbf{\textcolor{k8ssuccess}{Service Account Tokens}}
    \begin{itemize}
        \item Para aplicaciones que corren en el cluster
        \item JWT tokens automáticamente montados en pods
        \item Gestionados por el Service Account Admission Controller
    \end{itemize}
    
    \item \textbf{\textcolor{k8swarning}{Static Token File}}
    \begin{itemize}
        \item Tokens pre-compartidos en archivo
        \item No recomendado para producción
        \item Requiere reinicio del API server para cambios
    \end{itemize}
\end{enumerate}

\subsubsection{Creación de Certificados de Usuario}

\begin{commandbox}
\# Generar clave privada para usuario
openssl genrsa -out john.key 2048

\# Crear Certificate Signing Request
openssl req -new -key john.key -out john.csr -subj "/CN=john/O=developers"

\# Crear CertificateSigningRequest en Kubernetes
apiVersion: certificates.k8s.io/v1
kind: CertificateSigningRequest
metadata:
  name: john
spec:
  request: $(cat john.csr | base64 | tr -d '\n')
  signerName: kubernetes.io/kube-apiserver-client
  usages:
  - client auth

\# Aprobar el CSR
kubectl certificate approve john

\# Obtener el certificado
kubectl get csr john -o jsonpath='{.status.certificate}' | base64 -d > john.crt
\end{commandbox}

\subsubsection{Configuración de kubeconfig}

\begin{commandbox}
\# Configurar cluster
kubectl config set-cluster kubernetes --server=https://k8s-api:6443 \\
  --certificate-authority=ca.crt

\# Configurar usuario
kubectl config set-credentials john \\
  --client-certificate=john.crt \\
  --client-key=john.key

\# Configurar contexto
kubectl config set-context john@kubernetes \\
  --cluster=kubernetes \\
  --user=john \\
  --namespace=development

\# Usar contexto
kubectl config use-context john@kubernetes
\end{commandbox}

\subsection{Service Accounts}

Las Service Accounts proporcionan identidad para procesos que corren en pods.

\subsubsection{Gestión de Service Accounts}

\begin{commandbox}
\# Crear Service Account
kubectl create serviceaccount my-service-account

\# Ver Service Accounts
kubectl get serviceaccounts

\# Describir Service Account
kubectl describe serviceaccount my-service-account

\# Service Account declarativo
apiVersion: v1
kind: ServiceAccount
metadata:
  name: build-robot
  namespace: default
secrets:
- name: build-robot-secret
imagePullSecrets:
- name: my-registry-secret
\end{commandbox}

\subsubsection{Uso de Service Accounts en Pods}

\begin{commandbox}
apiVersion: v1
kind: Pod
metadata:
  name: my-pod
spec:
  serviceAccountName: my-service-account
  containers:
  - name: my-container
    image: nginx
    volumeMounts:
    - name: sa-token
      mountPath: /var/run/secrets/kubernetes.io/serviceaccount
      readOnly: true
  volumes:
  - name: sa-token
    projected:
      sources:
      - serviceAccountToken:
          path: token
          expirationSeconds: 3607
      - configMap:
          name: kube-root-ca.crt
          items:
          - key: ca.crt
            path: ca.crt
      - downwardAPI:
          items:
          - path: namespace
            fieldRef:
              fieldPath: metadata.namespace
\end{commandbox}

\subsection{Security Contexts}

Los Security Contexts definen configuraciones de seguridad para pods y contenedores.

\subsubsection{Security Context a nivel de Pod}

\begin{commandbox}
apiVersion: v1
kind: Pod
metadata:
  name: security-context-demo
spec:
  securityContext:
    runAsUser: 1000
    runAsGroup: 3000
    fsGroup: 2000
    seccompProfile:
      type: RuntimeDefault
  containers:
  - name: sec-ctx-demo
    image: busybox:1.28
    command: [ "sh", "-c", "sleep 1h" ]
    securityContext:
      allowPrivilegeEscalation: false
      readOnlyRootFilesystem: true
      runAsNonRoot: true
      capabilities:
        drop:
        - ALL
        add:
        - NET\_BIND\_SERVICE
\end{commandbox}

\subsubsection{Políticas de Seguridad Recomendadas}

\begin{tipbox}
\textbf{Mejores Prácticas de Seguridad:}
\begin{itemize}
    \item Usar usuarios no-root
    \item Deshabilitar escalation de privilegios
    \item Solo agregar capabilities necesarias
    \item Usar readOnlyRootFilesystem cuando sea posible
    \item Implementar seccomp profiles
\end{itemize}
\end{tipbox}

\subsection{Network Policies}

Las Network Policies actúan como firewalls para pods, controlando el tráfico de red.

\subsubsection{Componentes de Network Policy}

\begin{enumerate}
    \item \textbf{podSelector}: Selecciona pods a los que aplica la política
    \item \textbf{policyTypes}: Ingress, Egress, o ambos
    \item \textbf{ingress}: Reglas para tráfico entrante
    \item \textbf{egress}: Reglas para tráfico saliente
\end{enumerate}

\subsubsection{Ejemplo de Network Policy Completa}

\begin{commandbox}
apiVersion: networking.k8s.io/v1
kind: NetworkPolicy
metadata:
  name: test-network-policy
  namespace: default
spec:
  podSelector:
    matchLabels:
      role: db
  policyTypes:
  - Ingress
  - Egress
  ingress:
  - from:
    - ipBlock:
        cidr: 172.17.0.0/16
        except:
        - 172.17.1.0/24
    - namespaceSelector:
        matchLabels:
          name: web
    - podSelector:
        matchLabels:
          role: frontend
    ports:
    - protocol: TCP
      port: 6379
  egress:
  - to:
    - ipBlock:
        cidr: 10.0.0.0/24
    ports:
    - protocol: TCP
      port: 5978
\end{commandbox}

\subsubsection{Ejemplos Prácticos de Network Policies}

\begin{commandbox}
\# Denegar todo el tráfico
apiVersion: networking.k8s.io/v1
kind: NetworkPolicy
metadata:
  name: default-deny-all
spec:
  podSelector: {}
  policyTypes:
  - Ingress
  - Egress

\# Permitir solo tráfico DNS
apiVersion: networking.k8s.io/v1
kind: NetworkPolicy
metadata:
  name: allow-dns-access
spec:
  podSelector: {}
  policyTypes:
  - Egress
  egress:
  - to: []
    ports:
    - protocol: UDP
      port: 53
\end{commandbox}

\subsection{Pod Security Standards}

Los Pod Security Standards definen tres políticas para controlar aspectos de seguridad de pods.

\subsubsection{Niveles de Pod Security}

\begin{enumerate}
    \item \textbf{\textcolor{k8ssuccess}{Privileged}}: Sin restricciones (inseguro)
    \item \textbf{\textcolor{k8swarning}{Baseline}}: Previene escalaciones conocidas
    \item \textbf{\textcolor{k8serror}{Restricted}}: Fuertemente restrictivo (más seguro)
\end{enumerate}

\subsubsection{Implementación con Admission Controller}

\begin{commandbox}
\# Namespace con Pod Security Standards
apiVersion: v1
kind: Namespace
metadata:
  name: secure-namespace
  labels:
    pod-security.kubernetes.io/enforce: restricted
    pod-security.kubernetes.io/audit: restricted
    pod-security.kubernetes.io/warn: restricted
\end{commandbox}

\newpage

% ====== TROUBLESHOOTING & MONITORING ======
\section{Troubleshooting \& Monitoring}

\begin{infobox}
Esta sección cubre el 30\% del examen CKA y es crucial para la administración diaria de clusters Kubernetes.
\end{infobox}

\subsection{Troubleshooting de Clusters}

\subsubsection{Verificación del Estado del Cluster}

\begin{commandbox}
\# Estado general del cluster
kubectl cluster-info
kubectl get nodes -o wide
kubectl get pods --all-namespaces

\# Verificar componentes del control plane
kubectl get pods -n kube-system

\# Estado de componentes
kubectl get componentstatuses  # Deprecated pero útil

\# Información detallada de nodos
kubectl describe nodes
\end{commandbox}

\subsubsection{Troubleshooting de Nodos}

\begin{errorbox}
Problemas comunes: Nodos NotReady, recursos insuficientes, problemas de red, kubelet no funcional.
\end{errorbox}

\begin{commandbox}
\# Verificar estado de nodos
kubectl get nodes
kubectl describe node <node-name>

\# Logs del kubelet
sudo journalctl -u kubelet -f

\# Verificar recursos del sistema
top
free -h
df -h

\# Estado de contenedores
sudo crictl ps
sudo crictl pods

\# Reiniciar kubelet si es necesario
sudo systemctl restart kubelet
\end{commandbox}

\subsubsection{Troubleshooting de Pods}

\begin{commandbox}
\# Estados de pods problemáticos
kubectl get pods --field-selector=status.phase=Failed
kubectl get pods --field-selector=status.phase=Pending

\# Información detallada
kubectl describe pod <pod-name>
kubectl logs <pod-name> -c <container-name>
kubectl logs <pod-name> --previous

\# Debugging interactivo
kubectl exec -it <pod-name> -- /bin/bash
kubectl debug <pod-name> -it --image=busybox

\# Eventos del cluster
kubectl get events --sort-by='.lastTimestamp'
\end{commandbox}

\subsubsection{Problemas Comunes y Soluciones}

\begin{enumerate}
    \item \textbf{\textcolor{k8serror}{ImagePullBackOff}}
    \begin{itemize}
        \item Verificar nombre de imagen y tag
        \item Comprobar conectividad al registry
        \item Revisar image pull secrets
    \end{itemize}
    
    \item \textbf{\textcolor{k8serror}{CrashLoopBackOff}}
    \begin{itemize}
        \item Revisar logs del contenedor
        \item Verificar comando de inicio
        \item Comprobar health checks
    \end{itemize}
    
    \item \textbf{\textcolor{k8serror}{Pending}}
    \begin{itemize}
        \item Recursos insuficientes
        \item Node affinity no satisfecha
        \item Taints y tolerations
    \end{itemize}
\end{enumerate}

\subsection{Troubleshooting de Redes}

\subsubsection{Conectividad entre Pods}

\begin{commandbox}
\# Probar conectividad básica
kubectl run test-pod --image=busybox:1.28 --rm -it -- sh

\# Dentro del pod de prueba:
\# ping <pod-ip>
\# nslookup <service-name>
\# wget -qO- http://<service-name>:<port>

\# Verificar reglas iptables
sudo iptables -L -n -v

\# Estado de interfaces de red
ip addr show
ip route show
\end{commandbox}

\subsubsection{Troubleshooting de Services}

\begin{commandbox}
\# Verificar endpoints del service
kubectl get endpoints <service-name>
kubectl describe service <service-name>

\# Probar service desde dentro del cluster
kubectl run debug-pod --image=nicolaka/netshoot --rm -it -- sh

\# Verificar kube-proxy
kubectl get pods -n kube-system | grep kube-proxy
kubectl logs -n kube-system <kube-proxy-pod>

\# Verificar reglas de proxy
sudo iptables -t nat -L -n -v
\end{commandbox}

\subsection{Monitoring con Métricas}

\subsubsection{Metrics Server}

\begin{commandbox}
\# Instalar Metrics Server
kubectl apply -f \\
  https://github.com/kubernetes-sigs/metrics-server/releases/latest/download/components.yaml

\# Verificar instalación
kubectl get deployment metrics-server -n kube-system

\# Ver métricas de nodos
kubectl top nodes

\# Ver métricas de pods
kubectl top pods
kubectl top pods --all-namespaces
\end{commandbox}

\subsubsection{Comandos de Monitoreo Útiles}

\begin{commandbox}
\# Uso de recursos por namespace
kubectl top pods --sort-by=cpu -n <namespace>
kubectl top pods --sort-by=memory -n <namespace>

\# Métricas en tiempo real
watch kubectl top nodes
watch kubectl top pods --all-namespaces

\# Descripción de recursos con límites
kubectl describe limitrange -n <namespace>
kubectl describe resourcequota -n <namespace>
\end{commandbox}

\subsubsection{Troubleshooting de Performance}

\begin{tipbox}
\textbf{Herramientas de Análisis:}
\begin{itemize}
    \item kubectl top para métricas básicas
    \item kubectl describe para eventos y estado
    \item Logs de aplicación para errores específicos
    \item Herramientas del sistema (htop, iostat, netstat)
\end{itemize}
\end{tipbox}

\subsection{Backup y Disaster Recovery}

\subsubsection{Estrategias de Backup}

\begin{enumerate}
    \item \textbf{Backup de etcd} (más crítico)
    \item \textbf{Backup de configuraciones} (YAML manifests)
    \item \textbf{Backup de datos persistentes}
    \item \textbf{Backup de certificados y secrets}
\end{enumerate}

\subsubsection{Script de Backup Automatizado}

\begin{commandbox}
\#!/bin/bash
\# Automated backup script

DATE=$(date +%Y%m%d_%H%M%S)
BACKUP_DIR="/backup/$DATE"
mkdir -p $BACKUP_DIR

\# etcd backup
ETCDCTL\_API=3 etcdctl snapshot save $BACKUP_DIR/etcd-snapshot.db \\
  --endpoints=https://127.0.0.1:2379 \\
  --cacert=/etc/kubernetes/pki/etcd/ca.crt \\
  --cert=/etc/kubernetes/pki/etcd/server.crt \\
  --key=/etc/kubernetes/pki/etcd/server.key

\# Backup de manifests
kubectl get all --all-namespaces -o yaml > $BACKUP_DIR/all-resources.yaml

\# Backup de secrets
kubectl get secrets --all-namespaces -o yaml > $BACKUP_DIR/all-secrets.yaml

\# Backup de configmaps
kubectl get configmaps --all-namespaces -o yaml > $BACKUP_DIR/all-configmaps.yaml

echo "Backup completed in $BACKUP_DIR"
\end{commandbox}

\newpage

% ====== COMANDOS ESENCIALES Y CHEAT SHEET ======
\section{Comandos Esenciales y Cheat Sheet}

\subsection{Comandos Fundamentales}

\subsubsection{Gestión Básica}

\begin{commandbox}
\# Información del cluster
kubectl cluster-info
kubectl version
kubectl config view

\# Gestión de contextos
kubectl config get-contexts
kubectl config current-context
kubectl config use-context <context-name>

\# Namespaces
kubectl get namespaces
kubectl create namespace <name>
kubectl delete namespace <name>
\end{commandbox}

\subsubsection{Pods y Deployments}

\begin{commandbox}
\# Pods
kubectl run <pod-name> --image=<image> --port=<port>
kubectl get pods -o wide
kubectl describe pod <pod-name>
kubectl logs <pod-name> -f
kubectl exec -it <pod-name> -- /bin/bash
kubectl delete pod <pod-name>

\# Deployments
kubectl create deployment <name> --image=<image>
kubectl get deployments
kubectl scale deployment <name> --replicas=<number>
kubectl rollout status deployment/<name>
kubectl rollout history deployment/<name>
kubectl rollout undo deployment/<name>
\end{commandbox}

\subsubsection{Services y Networking}

\begin{commandbox}
\# Services
kubectl expose deployment <name> --port=<port> --type=<type>
kubectl get services
kubectl describe service <service-name>

\# Ingress
kubectl get ingress
kubectl describe ingress <ingress-name>

\# Network troubleshooting
kubectl run netshoot --image=nicolaka/netshoot --rm -it -- bash
kubectl port-forward pod/<pod-name> <local-port>:<pod-port>
\end{commandbox}

\subsection{Comandos de Troubleshooting}

\begin{commandbox}
\# Estado del cluster
kubectl get nodes
kubectl get pods --all-namespaces
kubectl get events --sort-by='.lastTimestamp'

\# Debugging de pods
kubectl describe pod <pod-name>
kubectl logs <pod-name> --previous
kubectl exec -it <pod-name> -- sh

\# Métricas
kubectl top nodes
kubectl top pods --all-namespaces

\# Debugging de nodos
kubectl describe node <node-name>
kubectl get pods --field-selector spec.nodeName=<node-name>
\end{commandbox}

\subsection{Comandos de Administración}

\begin{commandbox}
\# Drain y cordon de nodos
kubectl drain <node-name> --ignore-daemonsets --delete-emptydir-data
kubectl cordon <node-name>
kubectl uncordon <node-name>

\# Etiquetas y annotations
kubectl label nodes <node-name> <key>=<value>
kubectl annotate pods <pod-name> <key>=<value>

\# Secrets y ConfigMaps
kubectl create secret generic <name> --from-literal=<key>=<value>
kubectl create configmap <name> --from-file=<file>
\end{commandbox}

\newpage

% ====== EJEMPLOS PRÁCTICOS DEL EXAMEN CKA ======
\section{Ejemplos Prácticos del Examen CKA}

\begin{warningbox}
Esta sección incluye escenarios prácticos similares a los del examen CKA real. Practica estos ejercicios para prepararte adecuadamente.
\end{warningbox}

\subsection{Escenarios de Cluster Architecture}

\subsubsection{Escenario 1: Backup y Restore de etcd}

\textbf{Problema:} Realizar un backup de etcd y restaurarlo en caso de falla.

\begin{commandbox}
\# 1. Realizar backup de etcd
ETCDCTL\_API=3 etcdctl snapshot save /opt/etcd-backup.db \\
  --endpoints=https://127.0.0.1:2379 \\
  --cacert=/etc/kubernetes/pki/etcd/ca.crt \\
  --cert=/etc/kubernetes/pki/etcd/server.crt \\
  --key=/etc/kubernetes/pki/etcd/server.key

\# 2. Verificar el backup
ETCDCTL\_API=3 etcdctl snapshot status /opt/etcd-backup.db

\# 3. Restaurar (en caso necesario)
sudo systemctl stop kubelet
sudo systemctl stop etcd

ETCDCTL\_API=3 etcdctl snapshot restore /opt/etcd-backup.db \\
  --data-dir=/var/lib/etcd-restore

\# 4. Actualizar manifest de etcd y reiniciar
sudo vim /etc/kubernetes/manifests/etcd.yaml
sudo systemctl start kubelet
\end{commandbox}

\subsubsection{Escenario 2: Actualización de Cluster}

\textbf{Problema:} Actualizar un cluster de Kubernetes de v1.28 a v1.29.

\begin{commandbox}
\# 1. Verificar versión actual
kubectl get nodes

\# 2. Actualizar control plane
sudo apt update
sudo apt-mark unhold kubeadm
sudo apt install -y kubeadm=1.29.0-00
sudo apt-mark hold kubeadm

\# 3. Planificar actualización
sudo kubeadm upgrade plan

\# 4. Aplicar actualización
sudo kubeadm upgrade apply v1.29.0

\# 5. Actualizar kubelet y kubectl
kubectl drain control-plane --ignore-daemonsets
sudo apt-mark unhold kubelet kubectl
sudo apt install -y kubelet=1.29.0-00 kubectl=1.29.0-00
sudo apt-mark hold kubelet kubectl
sudo systemctl daemon-reload
sudo systemctl restart kubelet
kubectl uncordon control-plane
\end{commandbox}

\subsection{Escenarios de Networking}

\subsubsection{Escenario 3: Configurar Ingress}

\textbf{Problema:} Exponer una aplicación usando Ingress con path-based routing.

\begin{commandbox}
\# 1. Crear deployment
kubectl create deployment web-app --image=nginx:1.20
kubectl expose deployment web-app --port=80 --target-port=80

\# 2. Crear Ingress
apiVersion: networking.k8s.io/v1
kind: Ingress
metadata:
  name: web-app-ingress
  annotations:
    nginx.ingress.kubernetes.io/rewrite-target: /
spec:
  ingressClassName: nginx
  rules:
  - host: myapp.example.com
    http:
      paths:
      - path: /app
        pathType: Prefix
        backend:
          service:
            name: web-app
            port:
              number: 80

\# 3. Aplicar y verificar
kubectl apply -f ingress.yaml
kubectl get ingress
\end{commandbox}

\subsection{Escenarios de Security}

\subsubsection{Escenario 4: Configurar RBAC}

\textbf{Problema:} Crear un usuario que solo pueda leer pods en el namespace "development".

\begin{commandbox}
\# 1. Crear namespace
kubectl create namespace development

\# 2. Crear Role
apiVersion: rbac.authorization.k8s.io/v1
kind: Role
metadata:
  namespace: development
  name: pod-reader
rules:
- apiGroups: [""]
  resources: ["pods"]
  verbs: ["get", "watch", "list"]

\# 3. Crear RoleBinding
apiVersion: rbac.authorization.k8s.io/v1
kind: RoleBinding
metadata:
  name: read-pods
  namespace: development
subjects:
- kind: User
  name: dev-user
  apiGroup: rbac.authorization.k8s.io
roleRef:
  kind: Role
  name: pod-reader
  apiGroup: rbac.authorization.k8s.io

\# 4. Verificar permisos
kubectl auth can-i get pods --namespace=development --as=dev-user
\end{commandbox}

\subsection{Escenarios de Troubleshooting}

\subsubsection{Escenario 5: Pod no puede iniciar}

\textbf{Problema:} Un pod está en estado "Pending" y no puede programarse.

\begin{commandbox}
\# 1. Verificar estado del pod
kubectl get pods
kubectl describe pod <pod-name>

\# 2. Verificar eventos
kubectl get events --sort-by='.lastTimestamp'

\# 3. Verificar recursos de nodos
kubectl top nodes
kubectl describe nodes

\# 4. Verificar taints y tolerations
kubectl get nodes -o custom-columns=NAME:.metadata.name,TAINTS:.spec.taints

\# 5. Soluciones comunes:
\# - Agregar resources requests
\# - Agregar tolerations para taints
\# - Verificar nodeSelector/affinity
\end{commandbox}

\newpage

% ====== APÉNDICES ======
\section{Apéndices}

\subsection{Recursos Útiles}

\subsubsection{Documentación Oficial}

\begin{itemize}
    \item \href{https://kubernetes.io/docs/}{Documentación de Kubernetes}
    \item \href{https://kubernetes.io/docs/reference/kubectl/cheatsheet/}{Kubectl Cheat Sheet}
    \item \href{https://kubernetes.io/docs/concepts/}{Conceptos de Kubernetes}
    \item \href{https://kubernetes.io/docs/tasks/}{Tareas y Tutoriales}
\end{itemize}

\subsubsection{Herramientas de Práctica}

\begin{itemize}
    \item \href{https://killer.sh/}{killer.sh} - Simulador de examen CKA
    \item \href{https://kodekloud.com/}{KodeKloud} - Laboratorios interactivos
    \item \href{https://github.com/kelseyhightower/kubernetes-the-hard-way}{Kubernetes The Hard Way}
\end{itemize}

\subsection{Configuración de Laboratorio}

\subsubsection{Ambiente con Vagrant}

\begin{commandbox}
\# Vagrantfile para laboratorio K8s
Vagrant.configure("2") do |config|
  config.vm.box = "ubuntu/20.04"
  
  # Master node
  config.vm.define "master" do |master|
    master.vm.hostname = "k8s-master"
    master.vm.network "private_network", ip: "192.168.56.10"
    master.vm.provider "virtualbox" do |v|
      v.memory = 2048
      v.cpus = 2
    end
  end
  
  # Worker nodes
  (1..2).each do |i|
    config.vm.define "worker#{i}" do |worker|
      worker.vm.hostname = "k8s-worker#{i}"
      worker.vm.network "private_network", ip: "192.168.56.1#{i}"
      worker.vm.provider "virtualbox" do |v|
        v.memory = 1024
        v.cpus = 1
      end
    end
  end
end
\end{commandbox}

\subsubsection{Script de Instalación Automatizada}

\begin{commandbox}
\#!/bin/bash
\# install-k8s.sh - Script de instalación automatizada

set -e

\# Deshabilitar swap
sudo swapoff -a
sudo sed -i 's|/ swap |/# swap |g' /etc/fstab

\# Configurar módulos del kernel
cat <<EOF | sudo tee /etc/modules-load.d/k8s.conf
overlay
br\_netfilter
EOF

sudo modprobe overlay
sudo modprobe br\_netfilter

\# Configurar parámetros del kernel
cat <<EOF | sudo tee /etc/sysctl.d/k8s.conf
net.bridge.bridge-nf-call-iptables  = 1
net.bridge.bridge-nf-call-ip6tables = 1
net.ipv4.ip\_forward                 = 1
EOF

sudo sysctl --system

\# Instalar containerd
sudo apt-get update
sudo apt-get install -y ca-certificates curl gnupg lsb-release
curl -fsSL https://download.docker.com/linux/ubuntu/gpg | sudo gpg --dearmor -o /usr/share/keyrings/docker-archive-keyring.gpg

echo \\
  "deb [arch=$(dpkg --print-architecture) signed-by=/usr/share/keyrings/docker-archive-keyring.gpg] https://download.docker.com/linux/ubuntu \\
  $(lsb_release -cs) stable" | sudo tee /etc/apt/sources.list.d/docker.list > /dev/null

sudo apt-get update
sudo apt-get install -y containerd.io

\# Configurar containerd
sudo mkdir -p /etc/containerd
containerd config default | sudo tee /etc/containerd/config.toml
sudo sed -i 's/SystemdCgroup = false/SystemdCgroup = true/' /etc/containerd/config.toml
sudo systemctl restart containerd
sudo systemctl enable containerd

\# Instalar kubeadm, kubelet, kubectl
curl -fsSL https://packages.cloud.google.com/apt/doc/apt-key.gpg | sudo gpg --dearmor -o /etc/apt/keyrings/kubernetes-archive-keyring.gpg

echo "deb [signed-by=/etc/apt/keyrings/kubernetes-archive-keyring.gpg] https://apt.kubernetes.io/ kubernetes-xenial main" | sudo tee /etc/apt/sources.list.d/kubernetes.list

sudo apt-get update
sudo apt-get install -y kubelet kubeadm kubectl
sudo apt-mark hold kubelet kubeadm kubectl

echo "Instalación completada. Ejecutar 'kubeadm init' para inicializar el cluster."
\end{commandbox}

\subsection{Glosario de Términos}

\begin{itemize}
    \item \textbf{API Server}: Componente que expone la API de Kubernetes
    \item \textbf{CNI}: Container Network Interface - especificación para plugins de red
    \item \textbf{CRI}: Container Runtime Interface - especificación para runtimes
    \item \textbf{DaemonSet}: Asegura que todos los nodos ejecuten una copia de un pod
    \item \textbf{etcd}: Base de datos distribuida que almacena el estado del cluster
    \item \textbf{HPA}: Horizontal Pod Autoscaler - escala pods automáticamente
    \item \textbf{RBAC}: Role-Based Access Control - control de acceso basado en roles
    \item \textbf{Service Mesh}: Capa de infraestructura para comunicación entre servicios
    \item \textbf{StatefulSet}: Gestiona aplicaciones con estado
    \item \textbf{VPA}: Vertical Pod Autoscaler - ajusta recursos de pods automáticamente
\end{itemize}

\newpage

% ====== CONCLUSIÓN ======
\section{Conclusión}

\begin{infobox}
Esta guía completa de Kubernetes CKA 2024 cubre todos los aspectos esenciales para la administración de clusters Kubernetes y la preparación para el examen CKA.
\end{infobox}

\subsection{Puntos Clave para el Examen}

\begin{enumerate}
    \item \textbf{\textcolor{k8sblue}{Práctica Hands-on}}: El examen es 100\% práctico
    \item \textbf{\textcolor{k8ssuccess}{Tiempo de Gestión}}: Tienes 2 horas para 15-20 tareas
    \item \textbf{\textcolor{k8swarning}{Documentación}}: Puedes usar la documentación oficial
    \item \textbf{\textcolor{k8ssteelblue}{kubectl}}: Domina los comandos esenciales
    \item \textbf{\textcolor{k8serror}{Troubleshooting}}: Representa el 30\% del examen
\end{enumerate}

\subsection{Plan de Estudio Recomendado}

\begin{tipbox}
\textbf{Semana 1-2}: Arquitectura y conceptos básicos\\
\textbf{Semana 3-4}: Instalación y configuración\\
\textbf{Semana 5-6}: Networking y storage\\
\textbf{Semana 7-8}: Security y RBAC\\
\textbf{Semana 9-10}: Troubleshooting y práctica intensiva\\
\textbf{Semana 11-12}: Simulacros de examen
\end{tipbox}

\subsection{Comandos para Memorizar}

Los comandos más importantes que debes dominar:

\begin{commandbox}
kubectl get/describe/delete <resource>
kubectl create/apply -f <file>
kubectl logs/exec -it <pod>
kubectl drain/cordon/uncordon <node>
kubectl scale/rollout <deployment>
kubectl auth can-i <verb> <resource>
etcdctl snapshot save/restore
kubeadm init/join/upgrade
\end{commandbox}

\subsection{Última Recomendación}

\begin{warningbox}
La práctica constante en un ambiente real es la clave del éxito. Instala tu propio cluster, experimenta con todos los comandos y escenarios de esta guía. ¡Mucha suerte en tu examen CKA!
\end{warningbox}

\vfill
\begin{center}
\textcolor{k8sblue}{\Large\textbf{¡Feliz aprendizaje de Kubernetes!}}\\[0.5cm]
\textcolor{k8sdarkgray}{\faHeart\ Creado con amor para la comunidad de DevOps}
\end{center}
\end{document}
